% ===========================================================================
%  3 — Mapa dos épicos
% ===========================================================================
\chapter{Mapa dos épicos}
\label{cap:mapa}

São \textbf{14 épicos, 75 features e 229 histórias}. Os épicos seguem o
caminho de uso da plataforma (Figura~\ref{fig:caminho}): da entrada à
organização da rede, da sala de aula à missão, dos resultados às ferramentas,
as garantias que valem para tudo, a sala em que os estudantes dividem o
computador e, por fim, o que cada estudante demonstrou de cada habilidade da
BNCC.

\begin{destaque}
\textbf{Em 16 de outubro de 2026:} \textbf{155} histórias prontas
(\textasciitilde67\%), \textbf{57} parciais, \textbf{12} a fazer e \textbf{6}
arquivadas (as três do UnooTube, que saiu da trilha sem ser apagado, as perguntas do fim do Unoovista, as do resultado da Missão Tempestade, e a rodada
de desvio do Júri Estadual). Nove das
parciais são do agente de IA do professor (\id{F08.4} e \id{H10.1.4}),
construído e ainda não no ar, e oito são do Júri Estadual de Projetos
(\id{F06.14}, com o podcast dele e as imagens da mostra pronta), construído com
os quarenta estandes e com o teste de ponta a ponta à espera de ser atualizado
e corrido. Quatro são da Ajuda com atendente (\id{F01.6} e
\id{H10.1.5}), construída e em teste interno, desligada nas redes. Mais da
metade do que é parcial é a evidência pedagógica (\id{E14}) e as três missões
com régua, a versão nova de Quem mandou isso? (\id{F06.11}), A Seleção
(\id{F06.12}, com o podcast dela) e o Jornal da Escola (\id{H06.10.4}) ---
construídas e percorridas pelo teste de ponta a ponta; a primeira à espera de o
PO revisar os textos, a segunda à espera da revisão da voz e das imagens
geradas, da conferência dos níveis que a IA dá à legenda, da revisão dos
textos e da foto da Amara pelo PO, e a terceira à espera das imagens do banco padrão, da revisão
do critério do livro do 4º ano, da menção à IA dela nos Termos de Uso e da
primeira corrida verde do teste da aprovação das imagens (\id{H06.10.6}). O resto
é, em quase todos os casos, material que falta (áudio, foto, revisão
pedagógica) ou ajuste de tela. O que é a fazer concentra-se em recuperação de
senha, prazos de guarda de dados, acessibilidade, nas equipes do Unoovista e
nas missões que faltam. Oito itens --- de teste de carga a contrato --- bloqueiam a
entrada em produção (Capítulo~\ref{cap:backlog}).
\end{destaque}

\section{Os catorze épicos}

\begin{tabularx}{\linewidth}{@{}L{10mm}X L{12mm}L{9mm}L{9mm}L{9mm}L{17mm}@{}}
\toprule
\headrow \thd{ID} & \thd{Épico} & \thd{Feat.} & \thd{Pr.} & \thd{Parc.} & \thd{A f.} & \thd{Estado}\\
\midrule
\endhead
\id{E01} & Entrar e circular na plataforma (p.~\pageref{ep:e01}) & 6 & 10 & 3 & 2 & \parcial\\
\zebra \id{E02} & Organizar a rede (p.~\pageref{ep:e02}) & 4 & 7 & 0 & 1 & \parcial\\
\id{E03} & Cadastrar pessoas e turmas (p.~\pageref{ep:e03}) & 8 & 16 & 1 & 0 & \pronto\\
\zebra \id{E04} & Sala de aula: o professor abre a missão (p.~\pageref{ep:e04}) & 4 & 8 & 0 & 1 & \pronto\\
\id{E05} & A experiência do estudante na plataforma (p.~\pageref{ep:e05}) & 2 & 6 & 0 & 0 & \pronto\\
\zebra \id{E06} & As missões do Entre Redes (p.~\pageref{ep:e06}) & 12 & 43 & 28 & 1 & \pronto\\
\id{E07} & Acompanhar a aprendizagem e a rede (p.~\pageref{ep:e07}) & 3 & 6 & 1 & 1 & \parcial\\
\zebra \id{E08} & Ferramentas do professor (p.~\pageref{ep:e08}) & 4 & 10 & 8 & 0 & \pronto\\
\id{E09} & Inteligência artificial com limites (p.~\pageref{ep:e09}) & 4 & 5 & 0 & 2 & \parcial\\
\zebra \id{E10} & Painel da Unooria (p.~\pageref{ep:e10}) & 7 & 15 & 2 & 1 & \pronto\\
\id{E11} & Site público e captação (p.~\pageref{ep:e11}) & 3 & 3 & 1 & 0 & \parcial\\
\zebra \id{E12} & Confiança: privacidade, segurança e operação (p.~\pageref{ep:e12}) & 7 & 3 & 3 & 1 & \parcial\\
\id{E13} & Equipes na missão (p.~\pageref{ep:e13}) & 4 & 13 & 0 & 2 & \parcial\\
\zebra \id{E14} & Evidência pedagógica (p.~\pageref{ep:e14}) & 7 & 10 & 10 & 0 & \parcial\\
\midrule
 & \textbf{Total} & \textbf{75} & \textbf{155} & \textbf{57} & \textbf{12} & \\
\bottomrule
\end{tabularx}

\medskip
{\small\color{unooMuted}Pr. = prontas; Parc. = parciais; A f. = a fazer. As 6
histórias \arquivada{} (o UnooTube, \id{F06.2}, as perguntas do fim do Unoovista, \id{H06.6.8}, as do resultado da Missão Tempestade, \id{H06.8.5}, e a rodada de desvio do Júri,
\id{H06.14.8}) não entram nas colunas. O
estado do épico é \pronto{} quando o que falta é ajuste ou evolução e não
impede o uso; é \parcial{} quando falta uma capacidade que alguém vai sentir
falta.}

\section{Quem aparece em cada épico}

\begin{tabularx}{\linewidth}{@{}L{10mm}X *{7}{>{\centering\arraybackslash}p{9mm}}@{}}
\toprule
\headrow \thd{ID} & \thd{Épico} & \thd{Est.} & \thd{Prof.} & \thd{C.\,esc.} & \thd{C.\,reg.} & \thd{Adm.\,c.} & \thd{Uno.} & \thd{Vis.}\\
\midrule
\id{E01} & Entrar e circular & \yes & \yes & \yes & \yes & \yes & \yes & \no\\
\zebra \id{E02} & Organizar a rede & \no & \no & \yes & \yes & \yes & \yes & \no\\
\id{E03} & Pessoas e turmas & \no & \no & \yes & \yes & \yes & \yes & \no\\
\zebra \id{E04} & Sala de aula & \no & \yes & \yes & \no & \no & \no & \no\\
\id{E05} & Experiência do estudante & \yes & \no & \no & \no & \no & \no & \no\\
\zebra \id{E06} & Missões & \yes & \yes & \no & \no & \no & \no & \no\\
\id{E07} & Acompanhamento & \no & \yes & \no & \yes & \yes & \yes & \no\\
\zebra \id{E08} & Ferramentas do professor & \no & \yes & \no & \no & \no & \no & \no\\
\id{E09} & IA com limites & \yes & \yes & \no & \no & \no & \yes & \no\\
\zebra \id{E10} & Painel da Unooria & \no & \no & \no & \no & \yes & \yes & \no\\
\id{E11} & Site público & \no & \no & \no & \no & \no & \yes & \yes\\
\zebra \id{E12} & Confiança & \yes & \yes & \yes & \yes & \yes & \yes & \no\\
\id{E13} & Equipes na missão & \yes & \yes & \yes & \yes & \yes & \yes & \no\\
\zebra \id{E14} & Evidência pedagógica & \yes & \yes & \yes & \yes & \yes & \yes & \no\\
\bottomrule
\end{tabularx}

\medskip
{\small\color{unooMuted}Est. = Estudante; Prof. = Professor; C.\,esc. =
Coordenação da escola; C.\,reg. = Coordenação regional; Adm.\,c. =
Administrador do contrato; Uno. = Administrador Unooria; Vis. = Visitante. Em
\id{E10}, o Administrador do contrato só lê o histórico; em \id{E13}, as
coordenações e as administrações só consultam as equipes; em \id{E14}, elas só
leem, e o estudante responde sem ver o próprio resultado.}

\section{O que cada papel pode, em resumo}

\begin{tabularx}{\linewidth}{@{}X *{6}{>{\centering\arraybackslash}p{10.5mm}}@{}}
\toprule
\headrow \thd{Capacidade} & \thd{Est.} & \thd{Prof.} & \thd{C.\,esc.} & \thd{C.\,reg.} & \thd{Adm.\,c.} & \thd{Uno.}\\
\midrule
\endhead
Jogar as missões liberadas & \yes & \no & \no & \no & \no & \no\\
\zebra Liberar e encerrar missão na turma & \no & \yes & vê & vê & vê & vê\\
Montar as equipes de uma missão na turma & \no & \yes & vê & vê & vê & vê\\
\zebra Confirmar e ajustar a evidência por habilidade da BNCC & \no & \yes & vê & vê & vê & vê\\
Ler os resultados da turma e de cada estudante & o seu & \yes & \no & \no & \no & \no\\
\zebra Ler o agregado (sem identificar criança) & \no & \no & \no & \yes & \yes & \yes\\
Plano de aula, planejamento, bem-estar, agente de IA & \no & \yes & \no & \no & \no & \no\\
\zebra Cadastrar pessoas, importar, resetar senha & \no & \no & \yes & \yes & \yes & \yes\\
Turmas, enturmamento, ano letivo & \no & lê & \yes & \yes & \yes & \yes\\
\zebra Perfil portátil & \no & \no & \yes & \yes & \yes & \yes\\
Transferir estudante de escola & \no & \no & \no & \no & \yes & \yes\\
\zebra Regionais, escolas, vínculos em lote & \no & \no & \no & \no & \yes & \yes\\
Criar Coordenação regional & \no & \no & \no & \no & \yes & \yes\\
\zebra Criar Administrador do contrato & \no & \no & \no & \no & \no & \yes\\
Ligar e desligar recursos, versões, limites de IA & \no & \no & \no & \no & lê & \yes\\
\zebra Prospecção & \no & \no & \no & \no & \no & \yes\\
Tirar dúvidas com o atendente da Ajuda & \no & \yes & \yes & \yes & \yes & \yes\\
\bottomrule
\end{tabularx}

\medskip
{\small\color{unooMuted}Ausências deliberadas: o bem-estar do professor não
tem nenhuma visão de gestão; a gestão não tem acesso à linha de uma criança
nas missões --- a exceção declarada é a evidência pedagógica, em que ela lê o
nível de cada habilidade, nunca texto (\id{H14.4.4}, \id{F14.7}); o Administrador do
contrato lê a rede em agregado, nunca a sala de aula --- e, no painel de
habilidades, desce da rede ao estudante como as coordenações (\id{H14.7.5}); e o
estudante não vê a própria evidência pedagógica.}
